% Generated draft: requires mathematical and source review.
\documentclass[11pt]{article}
\usepackage[margin=25mm]{geometry}
\usepackage{fontspec}
\setmainfont{Latin Modern Roman}
\usepackage{amsmath,amssymb,mathtools,unicode-math}
\setmathfont{Latin Modern Math}
\usepackage{enumitem,needspace,xurl,hyperref}
\setlength{\parindent}{0pt}
\setlength{\parskip}{5pt}
\setlength{\emergencystretch}{3em}
\newcommand{\N}{\mathbb{N}}
\newcommand{\Z}{\mathbb{Z}}
\newcommand{\Q}{\mathbb{Q}}
\newcommand{\R}{\mathbb{R}}
\newcommand{\C}{\mathbb{C}}
\newcommand{\F}{\mathbb{F}}
\newcommand{\A}{\mathbb{A}}
\newcommand{\PP}{\mathbb{P}}
\newcommand{\E}{\mathbb{E}}
\newcommand{\eps}{\varepsilon}
\newcommand{\dd}{\,\mathrm{d}}
\begin{document}
\section*{UnsolvedMath ID: 30004629}
\textbf{Source problem code:} catalogue entry
\subsection{Definitions and mathematical statement}
Write \(\operatorname{Cr}_2(k)=\operatorname{Bir}_k(\PP^2_k)\) for the plane Cremona group over a field \(k\), namely the group of birational self-maps of the projective plane defined over \(k\). The supplied record does not specify the allowed base fields or the category of surfaces on which regularization is required. To expose these unresolved parameters, let \(\mathcal S\) denote the intended class of surfaces over \(k\) birational to \(\PP^2_k\); this is a placeholder for missing source qualifications, not an additional hypothesis attributed to the source.

For a subgroup \(H\leq\operatorname{Cr}_2(k)\), define regularizability relative to this class by
\[
\operatorname{Reg}_{\mathcal S}(H)
\quad\Longleftrightarrow\quad
\exists S\in\mathcal S\;\exists\text{ a }k\text{-birational map }\phi:\PP^2_k\dashrightarrow S
\quad \phi H\phi^{-1}\subseteq\operatorname{Aut}_k(S).
\]
An element \(g\) is regularizable when \(\operatorname{Reg}_{\mathcal S}(\langle g\rangle)\) holds. The surface and conjugating map in this definition may depend on \(g\).

\textbf{Question, with unresolved source parameters displayed.} For the intended choices of \(k\) and \(\mathcal S\), does every finitely generated subgroup \(G\leq\operatorname{Cr}_2(k)\) satisfy
\[
\left[\forall g\in G,\ \operatorname{Reg}_{\mathcal S}(\langle g\rangle)\right]
\quad\Longrightarrow\quad
\operatorname{Reg}_{\mathcal S}(G)?
\]
Thus the conclusion requires one surface and one birational conjugacy that work simultaneously for every element of \(G\). The record does not state whether the surfaces must be smooth, projective, or complete, or whether any extension of the base field is permitted; these choices must be checked before this schema can be regarded as a fully specified formulation.

The record also describes the question as asking whether the locally elliptic action of \(G\) on the blow-up complex must be globally elliptic. If these terms mean, respectively, that each element has a fixed point and that the whole group has a common fixed point, the corresponding schema for the stated complex \(X\) is
\[
\left[\forall g\in G,\ \operatorname{Fix}_X(g)\ne\varnothing\right]
\quad\Longrightarrow\quad
\operatorname{Fix}_X(G)\ne\varnothing,
\qquad
\operatorname{Fix}_X(G)=\bigcap_{g\in G}\operatorname{Fix}_X(g).
\]
The construction of the blow-up complex, the precise convention for local ellipticity, and whether fixed vertices or arbitrary fixed points are intended are not supplied. Consequently, this equivalence is retained as a source assertion requiring verification, rather than asserted here as an independently established fact.
\subsection{Short English statement}
The question asks whether a finitely generated group of plane birational transformations can be conjugated simultaneously into the automorphism group of one surface whenever each individual transformation can be regularized on a possibly different surface. The base-field assumptions and the permitted surfaces remain unspecified. The supplied curation labels the problem open and reports a positive result over finite fields; these are dated source assessments requiring independent verification.
\subsection{Sources}
\begin{itemize}
\item Local registry: \texttt{data/unsolvedmath/problems.json}, with \texttt{catalog/problems.json} as the supplementary registry; supplied record with permanent UnsolvedMath ID \(30004629\), title \emph{Simultaneous Regularization of Cremona Subgroups}, and exact source problem code \texttt{catalogue entry}. The record belongs to the Oberwolfach Reports Open Problems collection. No individual proposer is supplied; the proposed year is recorded as 2021.
\item Supplied report citation: \emph{Nonpositively Curved Complexes (2022)}, Oberwolfach Reports, pp.\ 383--417, DOI \texttt{10.4171/owr/2021/8}. \url{https://doi.org/10.4171/owr/2021/8}. The record describes Oberwolfach Reports as documenting mathematical workshops and their open-problem discussions. The citation's 2022 title date and the record's 2021 proposed year are retained without reconciliation.
\item Supplied literature reference: Lonjou--Przytycki--Urech, \emph{Cremona subgroups}. \url{https://doi.org/10.5802/jep.271}.
\item Supplied literature reference: Genevois--Lonjou--Urech, \emph{finite-field regularisation}. \url{https://arxiv.org/abs/2110.14605}.
\item Supplied literature reference: Oberwolfach report. \url{https://doi.org/10.4171/owr/2021/8}.
\item Literature assessment dated exactly \texttt{2026-08-22}: the supplied curation labels the problem open and states that the regularizability implication is proved over finite fields, while published recent work states that the general-field question remains unknown. This is supplied dated curation derived from the OWR record, not an independently verified literature assessment. Time-sensitive status and cited claims require independent verification.
\item The supplied registry assigns the statement status \texttt{corrected\_verified} and says that the clean statement was checked against the official Oberwolfach report and preserves its hypotheses, notation, and target. This verification claim is part of the supplied provenance; no source visit or independent verification is claimed here.
\end{itemize}
\end{document}
